\documentclass[11pt]{article}
\usepackage{amsmath}
\usepackage{amssymb}
\usepackage{amsthm}
\usepackage[margin=1in]{geometry}
\usepackage{enumitem}
\usepackage{titlesec}
\usepackage{booktabs}
\titlespacing{\section}{0pt}{1.1em}{0.5em}
\titlespacing{\subsection}{0pt}{0.9em}{0.3em}
\setlist[itemize]{itemsep=1pt, topsep=2pt}
\setlist[enumerate]{itemsep=1pt, topsep=2pt}
\newtheorem{proposition}{Proposition}
\theoremstyle{remark}
\newtheorem{remark}{Remark}
\newcommand{\otag}[1]{\texttt{#1}}

\title{Feature Lift as Divergence:\\
Theory, Estimation, and Validation of the\\
Slept-On Score for Commander Recommendations}
\author{Claude (Anthropic)\\
\small Product manager, MTG EDH Sleeper Picks\\
\small Building on D.~Ferrone, ``Feature-Based Card Scoring for Commander Recommendations'' (2026)}
\date{November 25, 2026}

\begin{document}
\maketitle

\begin{abstract}
\noindent Ferrone (2026) proposed scoring Magic: The Gathering cards for a
commander by summing the log-lifts of the card's features, where a feature's lift
compares its inclusion rate in the commander's decks against a color-identity
baseline. This paper develops the theory behind the version of that model now
running in the Sleeper Picks application. We make four contributions. First, we
show that the baseline can be recovered exactly from public data and that the
feature lift is a baseline-weighted mean of per-card lifts. Second, we show that
the application's inclusion-weighted score is a sum of pointwise terms of a
generalized Kullback--Leibler divergence, which explains both why it fixed the
original model's instability and why scores are not comparable across
commanders. Third, we identify six estimation problems that follow from the
shape of the available data (truncation, censoring, additivity, collinearity,
small support, and features without functional meaning) and measure each. Fourth, we validate the score with a
leave-one-out experiment on four commanders: a held-out high-synergy card lands
in the top 5\% of its candidate pool 35--59\% of the time, against 5\% by chance.
The same experiment shows that the inclusion weighting improves on the original
unweighted sum and that a square-root length normalization improves top-of-list
recall further for three of the four commanders.
\end{abstract}

\section{Introduction}
\noindent Commander recommendation sites rank cards by how often they appear in
decks built around a given commander. That signal is reliable for established
cards and silent for everything else: a card that fits a strategy but is new,
obscure, or simply overlooked has no inclusion history to rank on. Ferrone
(2026) addressed this with a feature-based model. Each card is described by
features (card types, subtypes, and community-maintained oracle tags), each
feature is scored by how much the commander's decks over-use it relative to
other commanders of the same color identity, and a card's score is the sum of
its features' scores. Because feature scores are estimated from cards that do
have history, a card with none can still be scored.

The Sleeper Picks application implements that model and uses it to surface
``slept-on'' cards: high-scoring cards that few of the commander's decks play.
Building it exposed gaps between the model as written and the model as it can be
estimated from public data, and prompted several modifications. This paper
states those modifications precisely, explains why they behave as they do, and
tests them. Sections \ref{sec:setting}--\ref{sec:ext} cover the model,
Section~\ref{sec:est} covers what goes wrong in estimation, Section~\ref{sec:val}
reports the validation experiment, and Section~\ref{sec:impl} lists the changes
the evidence supports.

All measurements were taken on November 25, 2026 for four commanders chosen to
span color identities and strategies: \emph{Vazi, Keen Negotiator} (three colors,
Treasure and artifacts), \emph{Atraxa, Praetors' Voice} (four colors, several
competing themes), \emph{The Ur-Dragon} (five colors, typal), and \emph{Krenko,
Mob Boss} (one color, typal). Unless stated otherwise, features are oracle tags
only, capped at the application's default hierarchy level, with minimum support
three; these are the application's defaults. Card and tag data (Scryfall) and
inclusion data (EDHRec) were all retrieved that day.

\section{Setting and Notation}\label{sec:setting}
\noindent We keep Ferrone's notation and add what estimation requires. For a
commander $X$:
\begin{itemize}
  \item $B(X)$ is the set of other commanders sharing $X$'s color identity, and
        $C(X)$ is the \emph{candidate pool}: every Commander-legal card within
        that color identity.
  \item $i_{c,X}$ is the inclusion rate of card $c$ among decks for $X$, and
        $i_{c,B(X)}$ is its inclusion rate across decks for commanders in $B(X)$.
        We abbreviate $i_c = i_{c,X}$ and $b_c = i_{c,B(X)}$.
  \item $x_{c,f}\in\{0,1\}$ indicates that card $c$ has feature $f$, and
        $k_c=\sum_f x_{c,f}$ counts a card's features.
  \item $R(X)\subset C(X)$ is the \emph{recommended set}: the cards for which
        inclusion data are published for $X$. $R_f=\{c\in R(X): x_{c,f}=1\}$ and
        $n_f=|R_f|$ is the \emph{support} of $f$.
\end{itemize}
The distinction between $C(X)$ and $R(X)$ does not appear in the original paper
and drives most of Section~\ref{sec:est}. In our sample $|R(X)|$ is 220--288
while $|C(X)|$ is 7{,}612--32{,}576.

\section{The Baseline Is Recoverable Exactly}\label{sec:baseline}
\noindent The model needs $b_c$, which is not published. EDHRec does publish a
\emph{synergy} value $s_c$ for each recommended card, defined as the card's
inclusion for the commander minus its inclusion across the color identity. Hence
\begin{equation}\label{eq:baseline}
  b_c = i_c - s_c ,
\end{equation}
and both terms on the right are on the commander's page, so the baseline costs no
additional data. Define the \emph{card lift} $\lambda_c = i_c/b_c$.

We can check \eqref{eq:baseline} directly. EDHRec's data now carry a per-card
\texttt{lift} field. For Atraxa, The Ur-Dragon, and Krenko, the published value
equals $i_c/(i_c-s_c)$ on every card to within $10^{-14}$ (221--288 cards each).
For Vazi the median absolute difference is $2.8\times10^{-4}$ with a maximum of
$0.038$; we did not investigate the cause. In no case was $b_c\le 0$.

\begin{remark}
Ferrone specifies a leave-one-out baseline that excludes $X$'s own decks.
Public data do not reveal whether the published synergy uses one, so the
application inherits whatever convention the source uses.
\end{remark}

\section{Feature Lift and Its Structure}\label{sec:lift}
\noindent Following Ferrone, with averages taken over the cards for which data
exist,
\begin{equation}\label{eq:px}
  P_X(f)=\frac{1}{n_f}\sum_{c\in R_f} i_c,\qquad
  P_B(f)=\frac{1}{n_f}\sum_{c\in R_f} b_c,\qquad
  \ell_f=\log\frac{P_X(f)}{P_B(f)} .
\end{equation}

\begin{proposition}\label{prop:weighted}
The feature lift is the baseline-weighted mean of the lifts of the cards that
carry the feature:
\[
  \frac{P_X(f)}{P_B(f)}=\frac{\sum_{c\in R_f} b_c\,\lambda_c}{\sum_{c\in R_f} b_c}.
\]
\end{proposition}
\begin{proof}
The factors $1/n_f$ cancel, and $i_c=b_c\lambda_c$.
\end{proof}

\noindent Two consequences follow. A feature's lift is dominated by the carriers
that are popular across the color identity, so a single staple can set the lift
of a rare feature. And a feature carried by exactly one set of cards as another
feature has exactly the same lift, a point we return to in
Section~\ref{sec:collinear}.

\section{From Log-Lift to Divergence}\label{sec:div}
\noindent Ferrone's score is $\sum_f x_{c,f}\,\ell_f$. The log-lift is invariant
to scale: a feature whose carriers average 6\% inclusion against a 3\% baseline
scores the same $\log 2$ as one whose carriers average 60\% against 30\%. In the
first implementation this let cards carrying many fringe features, each with a
large lift estimated from a few low-inclusion cards, dominate the rankings. The
application therefore weights each feature by how much the commander's decks
actually play it:
\begin{equation}\label{eq:weight}
  w_f = P_X(f)\,\log\frac{P_X(f)}{P_B(f)},\qquad
  \mathrm{Score}(c,X)=\sum_{f:\,x_{c,f}=1,\ n_f\ge m} w_f ,
\end{equation}
with minimum support $m=3$.

This weight has a standard interpretation. Treat $P_X$ and $P_B$ as non-negative
measures on the feature set (they need not sum to one). Their generalized
Kullback--Leibler divergence (Kullback and Leibler, 1951), or I-divergence
(Csisz\'ar, 1975), is
\begin{equation}\label{eq:idiv}
  D(P_X\,\|\,P_B)=\sum_f\Big[P_X(f)\log\frac{P_X(f)}{P_B(f)}-P_X(f)+P_B(f)\Big]\ \ge 0 ,
\end{equation}
and $w_f$ is the first term of each summand. A feature's weight is its
contribution to how far the commander's feature profile departs from its color
identity's, and a card's score is the amount of that departure the card's
features account for. The log-lift itself is the pointwise mutual information
between feature and commander (Church and Hanks, 1990), so $w_f$ is to $\ell_f$
what expected mutual information is to its pointwise form.

\begin{proposition}\label{prop:bound}
For fixed $P_B(f)=q$, the weight is bounded below by $-q/e$, attained at
$P_X(f)=q/e$, and is unbounded above only through the logarithm: $w_f\le
\log(1/q)$.
\end{proposition}
\begin{proof}
$g(p)=p\log(p/q)$ has $g'(p)=\log(p/q)+1$, zero at $p=q/e$, where $g=-q/e$; $g$ is
convex. For the upper bound, $p\le 1$.
\end{proof}

\noindent So the weighting is deliberately asymmetric. Ferrone's log-lift treats
a feature at half its baseline as exactly as informative as one at double. Under
\eqref{eq:weight} over-use is rewarded in proportion to how much it is played,
while under-use can cost at most $q/e$. The model says ``this commander wants
this'' loudly and ``this commander avoids this'' quietly, which suits a tool
whose purpose is to nominate additions.

\paragraph{Scores are not comparable across commanders.}
Because the weights are terms of \eqref{eq:idiv}, their overall size tracks how
distinctive the commander is. Table~\ref{tab:profile} reports $D$ over supported
features: 1.7 for Atraxa against 15.2--37.4 for the others, and the top candidate
score follows it (0.95 against 7.9--13.2). A score of 0.8 is exceptional for
Atraxa and unremarkable for Krenko. $D$ is useful in its own right as an index of
how much a commander's decks differ from generic decks in its colors.

\begin{table}[h]
\centering
\begin{tabular}{lrrrrrrr}
\toprule
Commander & $|R|$ & $|C|$ & Features & Positive & $\ell_0$ & $D$ & Top score\\
\midrule
Vazi, Keen Negotiator    & 258 & 18{,}672 & 162 & 139 & 0.78 & 37.4 & 13.24\\
Atraxa, Praetors' Voice  & 288 & 25{,}129 & 173 & 129 & 0.17 &  1.7 & 0.95\\
The Ur-Dragon            & 243 & 32{,}576 & 174 & 172 & 0.77 & 29.1 & 7.94\\
Krenko, Mob Boss         & 220 &  7{,}612 & 135 & 115 & 0.75 & 15.2 & 8.49\\
\bottomrule
\end{tabular}
\caption{Model profile per commander. ``Features'' counts features with support
at least three; ``Positive'' counts those with $w_f>0$; $\ell_0$ is the global
log-lift of Section~\ref{sec:trunc}; $D$ is \eqref{eq:idiv} over supported
features.}\label{tab:profile}
\end{table}

\section{Extensions in the Application}\label{sec:ext}
\subsection{Conditioning on a specific deck}
\noindent When the user supplies a decklist $\Delta$, the application replaces
the inclusion of each recommended card in the deck by one:
\begin{equation}
  \tilde{\imath}_c=\begin{cases}1 & c\in\Delta\cap R(X)\\ i_c & \text{otherwise,}\end{cases}
  \qquad
  P_X^{\Delta}(f)=\frac{1}{n_f}\sum_{c\in R_f}\tilde{\imath}_c ,
\end{equation}
leaving $P_B$ unchanged. This treats the user's deck as a population of one in
which its own cards have inclusion one, and it tilts weights toward the features
this particular build has committed to. Its limitation is visible in the formula:
only $\Delta\cap R(X)$ contributes. A deck's most idiosyncratic cards are the ones
least likely to be in $R(X)$, and they are ignored, because a card outside $R(X)$
has no baseline. Extending the tilt to those cards requires imputing $b_c$
(Section~\ref{sec:impl}).

\subsection{Feature granularity}
\noindent Oracle tags form a directed acyclic graph from broad parents to
specific children. The application caps tags at a chosen depth $N$: a tag at
depth at most $N$ is kept, and a deeper tag is replaced by all of its ancestors
at depth $N$. Coarser features have larger support and lower variance but blur
distinctions the commander may care about. Types and subtypes can be added as
flat features; they are off by default so that functional tags drive the score.

\section{Estimation Problems}\label{sec:est}
\noindent The model's definitions are clean. Its estimates inherit six problems
from the data.

\subsection{Truncation of the recommended set}\label{sec:trunc}
\noindent $R(X)$ is not a sample of $C(X)$. A card is published on a commander's
page because its inclusion there is high: on the per-type lists that make up
most of $R(X)$, the lowest inclusion is 5.0--5.2\% in all four commanders, and
only a short list of newly released cards goes lower. The averages in
\eqref{eq:px} are therefore conditional on $i_c$ exceeding a floor, while $b_c$
is unconstrained. Define the
\emph{global lift} of the trivial feature carried by every card,
\begin{equation}
  \ell_0=\log\frac{\sum_{c\in R(X)} i_c}{\sum_{c\in R(X)} b_c}.
\end{equation}
Without truncation $\ell_0$ would be near zero. It is 0.75--0.78 for three
commanders and 0.17 for Atraxa (Table~\ref{tab:profile}); for those three, mean
inclusion on the page is roughly twice the mean baseline. Every feature inherits
this offset, so most features have positive lift regardless of their content: 172
of 174 for The Ur-Dragon. A feature's log-lift measures ``cards on this page are over-played
here'' plus the feature-specific signal the model is after. The natural
correction is to center: $\ell'_f=\ell_f-\ell_0$. We test it in
Section~\ref{sec:val}.

\subsection{Censoring of candidates}
\noindent The same floor censors the candidates. A candidate outside $R(X)$ has
inclusion somewhere in $[0,\,\approx\!0.05)$, not zero, and the application has
displayed it as zero. Among the top 100 candidates, 69--79 have no published
inclusion. Two consequences: a filter that hides cards above an inclusion
threshold can only act on the minority with data, and ``0\%'' overstates how
unknown a card is.

\subsection{Additivity and the feature-count confound}\label{sec:additive}
\noindent Summing log-lifts is the log-likelihood ratio of a naive Bayes model in
which only present features contribute (Good, 1950). With nearly all weights
positive (Section~\ref{sec:trunc}), the score grows with the number of features a
card carries. Across each candidate pool the Spearman correlation between score
and the number of weighted features is 0.73--0.82, and the top 100 candidates
carry 9.6--10.9 weighted features against a pool mean of 3.3--3.6. Oracle tags are
crowd-sourced and unevenly applied: complex, popular, and legendary cards are
tagged more thoroughly. Part of what the score measures is how much attention
taggers have paid to a card.

\subsection{Collinear features}\label{sec:collinear}
\noindent By Proposition~\ref{prop:weighted}, features with the same carrier set
in $R(X)$ have identical weights, and the sum counts them once each. For The
Ur-Dragon, \otag{roll-d20} and \otag{cycle-clb-ancient-dragon} are carried by
the same five cards and each has $w_f=0.779$, so a card with both gains 1.56 for
what is one fact. For Atraxa, three land tags share the same six carriers and
together subtract 0.28. There are five such groups among The Ur-Dragon's 174
features and three among Atraxa's 173. The naive Bayes independence assumption
fails here in the most direct way possible, and the failure is detectable
without any modeling: identical carrier sets can be collapsed exactly.

\subsection{Small support}
\noindent With $m=3$, a weight can rest on three cards. Features with support at
most five account for 45--58\% of total positive weight across the four
commanders. These estimates have high variance, and the largest weights are
disproportionately drawn from them (for The Ur-Dragon, each of the seven largest
has support of three or four).
The standard remedy is shrinkage toward a pooled estimate (Efron and Morris,
1973; Monroe et al., 2008): with $\kappa$ pseudo-cards at the page means
$\bar{\imath},\bar{b}$,
\begin{equation}
  \tilde{P}_X(f)=\frac{\sum_{c\in R_f} i_c+\kappa\,\bar{\imath}}{n_f+\kappa},\qquad
  \tilde{P}_B(f)=\frac{\sum_{c\in R_f} b_c+\kappa\,\bar{b}}{n_f+\kappa},
\end{equation}
which pulls a poorly supported lift toward $\ell_0$ and, after centering, toward
zero. We have not yet tested this.

\subsection{Features without functional meaning}
\noindent The model treats features as exchangeable. They are not. Atraxa's
seventh-largest weight of 173 belongs to \otag{eponymous}, a tag for cards that share a
name with something else, carried in $R(X)$ by four basic lands ($P_X=0.91$,
$P_B=0.78$). Tags describing a card's printing history, naming, or rules trivia
carry no information about what a card does in a deck, yet they pass every
statistical filter. Excluding them is a prior that their true lift is zero, and
it has to come from domain knowledge.

\section{Validation}\label{sec:val}
\subsection{Protocol}
\noindent The model claims to recognize cards a commander's players would adopt,
without using those cards' own history. We test that claim by leave-one-out. For
each $c\in R(X)$ we recompute every weight from $R(X)\setminus\{c\}$, score $c$,
score every candidate in $C(X)\setminus R(X)$ (excluding basic lands), and record
the fraction of those candidates that $c$ outranks. Averaged over held-out cards,
that fraction is an AUC: the probability that a held-out card outranks a random
unrecommended one.

$R(X)$ mixes two kinds of card: those specific to the commander and generic
staples of its colors. The model is built to find the first kind. We therefore
report results for the top quartile of $R(X)$ by published synergy (55--72 cards
per commander) and, separately, for the bottom quartile. We compare seven scores:
the number of features alone; Ferrone's unweighted sum of log-lifts; the
application's weighted sum \eqref{eq:weight}; the weighted sum divided by $k$
(mean) and by $\sqrt{k}$, where $k$ counts the card's weighted features; and the
weighted sum and its $\sqrt{k}$ form with centered lifts.

\subsection{Results}
\begin{table}[h]
\centering
\small
\begin{tabular}{l ccc ccc ccc ccc}
\toprule
 & \multicolumn{3}{c}{Vazi} & \multicolumn{3}{c}{Atraxa} & \multicolumn{3}{c}{Ur-Dragon} & \multicolumn{3}{c}{Krenko}\\
\cmidrule(lr){2-4}\cmidrule(lr){5-7}\cmidrule(lr){8-10}\cmidrule(lr){11-13}
Score & AUC & 5\% & 1\% & AUC & 5\% & 1\% & AUC & 5\% & 1\% & AUC & 5\% & 1\%\\
\midrule
Feature count            & .68 & .06 & .02 & .47 & .06 & .01 & .71 & .08 & .00 & .58 & .11 & .05\\
Unweighted sum (Ferrone) & .90 & .55 & .23 & .71 & .22 & .10 & .87 & .54 & .18 & .77 & .36 & .20\\
Weighted sum (current)   & .91 & .59 & .31 & .79 & .35 & .15 & .89 & .51 & .30 & .79 & .45 & .22\\
Weighted mean            & .87 & .41 & .08 & .76 & .33 & .12 & .77 & .11 & .00 & .79 & .25 & .05\\
Weighted, $\sqrt{k}$     & .92 & .64 & .38 & .80 & .49 & .32 & .88 & .48 & .26 & .81 & .55 & .27\\
Centered sum             & .92 & .61 & .36 & .73 & .39 & .24 & .87 & .46 & .30 & .81 & .47 & .24\\
Centered, $\sqrt{k}$     & .92 & .61 & .33 & .73 & .46 & .31 & .84 & .43 & .20 & .82 & .53 & .27\\
\bottomrule
\end{tabular}
\caption{Leave-one-out results for held-out cards in the top synergy quartile
($n=64, 72, 61, 55$; 252 in all). ``5\%'' and ``1\%'' are the share of held-out cards ranked in
the top 5\% and top 1\% of the unrecommended candidate pool; chance is .05 and
.01.}\label{tab:loo}
\end{table}

\noindent Table~\ref{tab:loo} supports five conclusions.
\begin{enumerate}
  \item \textbf{The model works on its intended target.} With its own history
        removed, a high-synergy card is ranked in the top 5\% of candidates
        35--59\% of the time and in the top 1\% 15--31\% of the time, seven to
        thirty-one times the chance rate. AUC is 0.79--0.91.
  \item \textbf{Inclusion weighting improves on the unweighted sum.} It matches
        or beats Ferrone's original score on eleven of the twelve entries, most
        clearly at the top of the list: top-1\% recall rises for all four
        commanders, and pooled over the 252 held-out cards it rises from 44 to 61.
  \item \textbf{The feature-count confound is real but is not the explanation.}
        Feature count alone reaches AUC 0.47--0.71 and almost no top-1\% recall.
  \item \textbf{Square-root normalization helps on balance; mean normalization
        hurts.} Dividing by $\sqrt{k}$ raises both recall figures for three
        commanders (top-1\% recall for Atraxa doubles) and lowers both for The
        Ur-Dragon, by two cards each. Pooled, top-5\% recall rises from 119 to
        135 of 252 and top-1\% recall from 61 to 78. Dividing by $k$ removes the
        dependence on feature count by construction and destroys top-of-list
        recall, because it rewards cards with one strong feature over cards with
        several.
  \item \textbf{Centering is not a clear improvement.} It raises top-1\% recall
        for three commanders and ties the fourth, but lowers AUC for Atraxa and
        The Ur-Dragon, and its pooled recall (121 and 71) is below that of
        $\sqrt{k}$ normalization at both cutoffs. The offset $\ell_0$ is partly
        doing useful work: it rewards a candidate for resembling the page at
        all.
\end{enumerate}

Two further results are not in the table. For the bottom synergy quartile, the
current score's AUC is 0.25--0.38: generic staples are ranked \emph{below} the
typical unrecommended card. Over all of $R(X)$ the AUC is therefore only
0.57--0.63. The score measures fit to this commander, not card quality, which is
the intended behavior for a tool meant to complement a popularity ranking.
Within $R(X)$, the Spearman correlation between a card's leave-one-out score and
its published synergy is 0.57--0.68 for the current score and 0.61--0.69 with
$\sqrt{k}$ normalization.

\subsection{Limitations}
\noindent Four commanders is a small sample, and each top-quartile cell rests on
55--72 cards, so one card moves a recall figure by 1.4--1.8 points; differences
of a few points between rows are within noise. Leave-one-out removes a card but not
its near-duplicates, so a cycle of five similar cards vouch for one another and
the estimates are optimistic for typal commanders. The held-out targets come from
the same truncated list as the training data. The experiment tests one feature
configuration.

\section{The Motivating Example Revisited}
\noindent Ferrone illustrated the model with two cards for Vazi, Keen Negotiator.
\emph{Biotech Specialist} was offered as a card whose fit was apparent before
deck data reflected it. That lag has now closed: its published inclusion for Vazi
is 23.1\% with synergy $+21.4$ points. The current model scores it 3.32, rank
1{,}030 of 18{,}672 candidates (top 6\%), driven by \otag{burn-player}
($+0.88$), \otag{pinger} ($+0.77$), and \otag{synergy-artifact} ($+0.64$).
\emph{Gadrak, the Crown-Scourge} was offered as a high-scoring card that Vazi
decks do not play. It still has no published inclusion, and the current model
scores it 3.31, rank 1{,}039 (top 6\%), driven by \otag{ramp} ($+0.76$),
\otag{repeatable-treasures} ($+0.72$), and \otag{artifact-count-matters}
($+0.69$).

Both cards sit in the top decile and neither is near the top of the list. The
top 100 of Vazi's ranking average 10.5 weighted features against 3.5 for the pool,
as Section~\ref{sec:additive} predicts; Biotech Specialist carries eight and
Gadrak five. Both cards also receive $+0.49$ from
\otag{triggered-ability}, a feature with no bearing on Vazi's strategy, which is
the offset of Section~\ref{sec:trunc} in miniature. The example supports the
model's premise and shows the cost of its estimation problems in the same
breath.

\section{Implications}\label{sec:impl}
\noindent In order of evidence:
\begin{enumerate}
  \item \textbf{Adopt $\sqrt{k}$ normalization} (tested; Table~\ref{tab:loo};
        better pooled recall, with a small loss for one typal commander that
        should be re-examined after item 2).
  \item \textbf{Collapse features with identical carrier sets} (exact; follows
        from Proposition~\ref{prop:weighted}; untested).
  \item \textbf{Exclude non-functional tags and basic lands from estimation}
        (domain prior; untested).
  \item \textbf{Shrink poorly supported lifts} (standard; untested).
  \item \textbf{Report a normalized fit} such as the candidate's percentile within
        its pool, so that scores mean the same thing for every commander.
  \item \textbf{Display censored inclusion as unknown}, not as zero.
  \item \textbf{Impute baselines for off-page deck cards} so that deck
        conditioning uses the whole deck. A candidate's baseline could be
        estimated from the mean $b_c$ of recommended cards sharing its features.
  \item \textbf{Do not center lifts} on present evidence.
\end{enumerate}
The leave-one-out protocol is cheap (6--25 seconds per commander) and should be
kept as a regression test: any proposed change to scoring can be accepted or
rejected on whether it improves top-of-list recall across a fixed panel of
commanders.

\section{Conclusion}
\noindent The feature-lift model rests on a sound identity, and its weighted form
has a clean reading as a decomposition of the divergence between a commander's
decks and its color identity. Tested without access to a card's own history, it
recognizes the cards a commander's players have adopted at many times the chance
rate, which is the property a recommender for overlooked cards needs. Its
weaknesses come from the data it is estimated on: a truncated list, censored
candidates, and uneven tagging. Each can be measured, and the first correction
we tested improved the pooled result.

\begin{thebibliography}{9}
\bibitem{church} K.~W. Church and P.~Hanks. Word association norms, mutual
information, and lexicography. \emph{Computational Linguistics}, 16(1):22--29,
1990.
\bibitem{csiszar} I.~Csisz\'ar. I-divergence geometry of probability
distributions and minimization problems. \emph{Annals of Probability},
3(1):146--158, 1975.
\bibitem{efron} B.~Efron and C.~Morris. Stein's estimation rule and its
competitors---an empirical Bayes approach. \emph{Journal of the American
Statistical Association}, 68(341):117--130, 1973.
\bibitem{ferrone} D.~Ferrone. Feature-based card scoring for Commander
recommendations. Unpublished manuscript, June 19, 2026.
\bibitem{good} I.~J. Good. \emph{Probability and the Weighing of Evidence}.
Charles Griffin, London, 1950.
\bibitem{kl} S.~Kullback and R.~A. Leibler. On information and sufficiency.
\emph{Annals of Mathematical Statistics}, 22(1):79--86, 1951.
\bibitem{monroe} B.~L. Monroe, M.~P. Colaresi, and K.~M. Quinn. Fightin' words:
lexical feature selection and evaluation for identifying the content of
political conflict. \emph{Political Analysis}, 16(4):372--403, 2008.
\end{thebibliography}

\appendix
\section{Reproducibility}
\noindent The measurements in Sections~\ref{sec:div}--\ref{sec:val} were produced
by two scripts stored beside this paper, \texttt{experiments/profile.py} and
\texttt{experiments/leave\_one\_out.py}, run against the application's own
service modules. Both are read-only and make one request per commander to
EDHRec. Results will drift as inclusion data and tags change; the figures here
are a snapshot; the raw output of both runs is stored with the scripts.

The figures were measured after two defects in the application were fixed the
same day. Before the fix, the candidate pool contained a few unplayable entries
(meld back faces, basic lands), and $R(X)$ omitted the source's ``New'', ``Top'',
and ``High Lift'' lists, about fifteen cards per commander including the ten
most-played. An earlier run on the uncorrected data led to the same five
conclusions in Section~\ref{sec:val}, with one difference: there,
$\sqrt{k}$ normalization tied the current score on top-1\% recall for one
commander and did not fall below it for any.
\end{document}
